\subsection{Function Notation} 
When we have a function, we have an alternate notation for it.
\begin{displaybox}[Function Notation]
If we have a relation $f$ that is a function, we can write
\[
f(x)
\]
read ``$f$ of $x$,'' to mean
\[
\mbox{the value of $y$ such that $(x,y)$ is an element of $f$}
\]
\end{displaybox}
\begin{Note} We usually use capital letters from the last third of the alphabet (\makeblanksmall, \makeblanksmall, \makeblanksmall, etc) for relations, but lower-case letters from the first third of the alphabet (\makeblanksmall, \makeblanksmall, \makeblanksmall, \makeblanksmall) for functions. These are conventions, not rules.
\end{Note}
\begin{Example}
Let $f$ be the relation
\[
f=\left\{(0,1),(1,-2),(2,3),(3,-4)\right\}
\]
Evaluate each of the following:
\begin{itemize}
\item $f(1)=\makeblanksmall$, since \blankpair is an element of $f$.
\item $f(3)=\makeblanksmall$, since \blankpair is an element of $f$.
\item $f(-2)$ is \makeblank[1in], since there $-2$ is not in the \makeblank[1in] of $f$.
\end{itemize}
\end{Example}
We can do the same thing when the function is given as a graph.
\begin{Example}The function $g$ is graphed below.
\begin{icpic}[x=0.5\linespace,y=0.5\linespace]
\setgraphsquare{5}
\GraphSmall
\draw[graphend,<-] (-5,-2) -- (-1,2);
\draw[graphend,->] (2,-4) -- (5,2);
\filldraw[dotfull] (-1,2) circle (0.2);
\filldraw[dotempty] (2,-4) circle (0.2);
\draw[graphend] (5,5) node[anchor=north east]{$g$};
\end{icpic}
\begin{itemize}
\item $f(-2)=\makeblanksmall$ because \blankpair is on the graph.
\item $f(3)=\makeblanksmall$ because \blankpair is on the graph.
\item $f(1)$ is \makeblank[1in], since $1$ is not in the \makeblank[1in] of $g$.
\end{itemize}
\end{Example}